\section{The Mise en Place Methodology}
\label{sec:methodology}

We propose \emph{mise en place} as a structured preparation methodology for agentic coding.
The term, borrowed from professional culinary practice, describes the discipline of
arranging every ingredient and tool before cooking begins. In a professional kitchen,
the speed of service is a direct function of the thoroughness of preparation:
chefs who invest time in mise en place execute complex dishes fluidly, while
those who skip preparation spend service time searching, measuring, and
recovering from avoidable errors. We observe an analogous dynamic in
AI-assisted software development:
when practitioners invest in deliberate preparation---externalizing domain knowledge,
producing detailed specifications, and decomposing work into structured task
records---the resulting agent execution is faster, more aligned with intent, and
requires less corrective iteration.

The methodology formalizes preparation into three sequential phases, each
producing defined artifacts that feed into the next. Following the backward design
principle~\cite{wiggins1998understanding}, the phases are completed before any
implementation begins. Figure~\ref{fig:methodology}
illustrates the three phases and their relationship to agent execution.

\begin{figure}[t]
  \centering
  \begin{tikzpicture}[
    >=Stealth,
    phase/.style={draw, rounded corners, fill=black!8, minimum width=2.1cm,
      minimum height=0.7cm, font=\footnotesize, align=center, text width=2cm},
    agent/.style={draw, rounded corners, fill=black!4, minimum width=1.2cm,
      minimum height=0.55cm, font=\scriptsize, align=center},
    label/.style={font=\scriptsize\itshape, text=black!60},
    arw/.style={->, thick},
  ]
    % Preparation phases
    \node[phase] (p1) {Contextual\\Grounding};
    \node[phase, right=0.35cm of p1] (p2) {Collaborative\\Specification};
    \node[phase, right=0.35cm of p2] (p3) {Task\\Decomposition};

    % Arrows between phases
    \draw[arw] (p1) -- (p2);
    \draw[arw] (p2) -- (p3);

    % Parallel agents
    \node[agent, below=0.9cm of p2, xshift=-0.7cm] (a1) {Agent 1};
    \node[agent, right=0.2cm of a1] (a2) {Agent 2};
    \node[agent, right=0.2cm of a2] (a3) {Agent $n$};

    % Arrow from decomposition to agents
    \draw[arw] (p3.south) -- ++(0,-0.35) -| (a2.north);

    % Testing node
    \node[phase, below=0.9cm of a2, minimum width=2.5cm] (test) {Testing \&\\Polish};

    % Arrows from agents to testing
    \draw[arw] (a1.south) -- ++(0,-0.25) -| ([xshift=-0.3cm]test.north);
    \draw[arw] (a2.south) -- (test.north);
    \draw[arw] (a3.south) -- ++(0,-0.25) -| ([xshift=0.3cm]test.north);

    % Feedback loop
    \draw[arw, dashed] (test.east) -- ++(0.35,0) |- (p3.east);

    % Labels
    \node[label, above=0.08cm of p2] {Preparation};
    \node[label, left=0.05cm of a1] {Execution};

  \end{tikzpicture}
  \caption{The mise en place methodology: three preparation phases produce structured context that enables parallel agent execution with minimal coordination overhead. Dashed arrow indicates a feedback loop for defects.}
  \label{fig:methodology}
\end{figure}

\subsection{Phase 1: Contextual Grounding}

The first phase externalizes domain expertise and tacit knowledge into structured
documents that agents can consume. Its purpose is to capture context that agents
cannot generate independently---what Polanyi~\cite{polanyi1966tacit} characterized
as tacit knowledge, the understanding that practitioners possess but struggle
to articulate. In agentic coding, this includes institutional context,
domain-specific constraints, competitive landscape awareness, design philosophy,
and the practitioner's accumulated judgment about what constitutes quality in
the target domain.

The artifacts produced are briefing documents: markdown files encoding domain
knowledge, competitive analysis, architectural preferences, and design philosophy.
These documents serve a dual purpose. First, they force the practitioner to
articulate knowledge that would otherwise remain implicit. Second, they create a
persistent context layer that agents reference throughout implementation, reducing
the need for repeated clarification. Following backward
design~\cite{wiggins1998understanding}, contextual grounding begins with outcomes
(the ``why'') rather than technical features (the ``what''), establishing alignment
criteria that propagate through subsequent phases.

\subsection{Phase 2: Collaborative Specification}

The second phase produces detailed design artifacts through structured
human-agent dialogue. The practitioner describes intent; the agent proposes
implementation details; the practitioner accepts, rejects, or modifies the
proposal. This iterative refinement continues until both parties converge on a
specification detailed enough for implementation without further clarification.

The key mechanism is the encoding of value judgments as specification constraints.
What to \emph{exclude} is as important as what to include. Scope
management---the deliberate decision to constrain features, simplify interactions,
or defer capabilities---represents a class of judgment that agents cannot make
independently. The output is a comprehensive design document that specifies
screens, interactions, data flows, and quality standards, capturing not just the
\emph{what} but the \emph{why} behind each decision. This distinction separates
collaborative specification from conventional requirements gathering: the
specification encodes rationale, enabling agents to make aligned micro-decisions
during implementation without escalating to the practitioner.

\subsection{Phase 3: Task Decomposition}

The final preparatory phase converts the specification into structured,
dependency-aware task records. We use
Beads~\cite{yegge2025beads, emanuel2025beadsrust}---lightweight, Git-backed
JSON records with priorities, dependencies, and acceptance criteria---though
the principle applies to any structured task tracking system.

Fine-grained decomposition enables parallel agent execution without coordination
overhead. When each task is independently executable---with clear boundaries,
defined inputs and outputs, and explicit dependencies---multiple agents can work
simultaneously. The coordination burden shifts from runtime (agents negotiating)
to preparation time (the practitioner defining task boundaries), where human
judgment about system architecture is most valuable. The output of this phase is a dependency graph where each node represents an
independently executable unit of work. The graph structure enables the
practitioner to identify the critical path, allocate agents to parallelizable
tasks, and monitor progress through structured completion signals. It also
provides a natural mechanism for prioritization: core features are
executed first, while stretch goals are explicitly deferred, preventing scope
creep during implementation.

The three phases of mise en place draw individually on established work---tacit
knowledge externalization~\cite{polanyi1958personal}, the RPI
methodology~\cite{horthy2025rpi}, and iterative task
delegation~\cite{huntley2025ralph}. The methodology's contribution is their
integration into a sequential, phase-gated process with defined artifacts and
theoretical grounding in backward design: the insistence that all three phases
complete before implementation begins.
